\oldpage[IV-3]{145}
et comme $M_0$ est par hypothèse un $A_0$-module \emph{plat}, le noyau de
$f_\lambda$ est égal à $M_0\otimes_{A_0}R_\lambda$, où $R_\lambda$ est le
noyau de
$\operatorname{gr}(u_\lambda):\operatorname{gr}(A)
\to\operatorname{gr}_{\mathfrak J}(B_\lambda)$. Tout revient donc à prouver
que
$\bigcap_{\lambda\in L}(M_0\otimes_{A_0}R_\lambda)=0$.
Or, par définition des $\mathfrak J_n$, l'intersection des noyaux des
homomorphismes
$\mathfrak J_nA/\mathfrak J_{n+1}A
\longrightarrow\mathfrak J^nB_\lambda/\mathfrak J^{n+1}B_\lambda$,
lorsque $\lambda$ parcourt $L$, est réduite à $0$, autrement dit
$\bigcap_{\lambda\in L}R_\lambda=0$. Cela étant, supposons d'abord que $M_0$
soit un $A_0$-module libre~; prenant une base de $M_0$, on voit aussitôt que
l'on a
$M_0\otimes_{A_0}(\bigcap_{\lambda\in L}R_\lambda)
=\bigcap_{\lambda\in L}(M_0\otimes_{A_0}R_\lambda)$,
d'où la proposition dans ce cas. Lorsque $L$ est fini, la formule précédente
est encore vraie sous la seule hypothèse que $M_0$ est un $A_0$-module plat
$(\mathbf 0_{\mathrm I},6.1.3)$, ce qui termine la démonstration.
\end{proof}

\begin{remark}[11.4.2]
\label{IV.11.4.2-fr}
La conclusion de (11.4.1) peut être inexacte si $L$ est infini et si l'on
suppose seulement que $M\otimes_A(A/\mathfrak J)$ est un
$(A/\mathfrak J)$-module plat. Par exemple, soient $V$ un anneau de valuation
discrète, $K$ son corps des fractions, et soit $A=V[T]/(T^2)$ ($T$
indéterminée)~; prenons pour $\mathfrak J$ l'image de $(T)$ dans $A$, de sorte
que $A/\mathfrak J=V$, et prenons $M=K$, qui est un
$(A/\mathfrak J)$-module, donc égal à $M\otimes_A(A/\mathfrak J)$~; en outre
$M$ est un $(A/\mathfrak J)$-module plat, mais non un $A$-module plat, car
puisque $\mathfrak J$ est nilpotent, il résulterait de
$(\mathbf 0_{\mathrm{III}},10.1.2)$ que $M$ serait un $A$-module libre, ce qui
est absurde puisque $\mathfrak JK=0$. Considérons d'autre part l'idéal maximal
$\mathfrak m$ de l'anneau local noethérien $A$~; on a $\mathfrak mK=K$, donc
$M\otimes_A(A/\mathfrak m^n)=0$ quel que soit l'entier $n$~; les
$(A/\mathfrak m^n)$-modules $M\otimes_A(A/\mathfrak m^n)$ sont donc plats pour
tout $n$, et l'intersection des $\mathfrak m^n$ est réduite à $0$.
\end{remark}

\begin{corollary}[11.4.3]
\label{IV.11.4.3-fr}
Soient $A$ un anneau semi-local dont le radical $\mathfrak J$ est nilpotent
(par exemple un anneau artinien), $u_\lambda:A\to B_\lambda$ $(\lambda\in L)$
une famille d'homomorphismes d'anneaux telle que l'intersection des noyaux des
$u_\lambda$ soit réduite à $0$. Pour qu'un $A$-module $M$ soit plat, il faut et
il suffit que pour tout $\lambda\in L$, $M\otimes_A B_\lambda$ soit un
$B_\lambda$-module plat.
\end{corollary}

\begin{proof}
Comme $A/\mathfrak J$ est composé direct d'un nombre fini de corps (Bourbaki,
\emph{Alg. comm.}, chap.~II, \S~3, n\textsuperscript{o}~5, prop.~16) et que $\mathfrak J$
est nilpotent, $A$ est composé direct d'un nombre fini d'anneaux \emph{locaux} $A_i$
dont le radical est nilpotent (\emph{loc. cit.}, \S~4, n\textsuperscript{o}~3, cor.~1 de
la prop.~15) et $M$ est par suite somme directe de $A_i$-modules $M_i$, chaque
$M_i$ étant annulé par les $A_j$ d'indice $j\neq i$~; pour que $M$ soit un
$A$-module plat, il faut et il suffit que chacun des $M_i$ soit un $A_i$-module
plat~; d'ailleurs l'intersection des noyaux des homomorphismes
$A_i\longrightarrow A\xrightarrow{u_\lambda}B_\lambda$
est réduite à $0$, et $M_i\otimes_{A_i}B_\lambda$ est facteur direct de
$M\otimes_A B_\lambda$. On peut donc se borner au cas où $A$ est en outre
\emph{local}. Alors $A/\mathfrak J$ est un corps, donc
$M\otimes_A(A/\mathfrak J)$ est un $(A/\mathfrak J)$-module libre, et il
suffit de voir que pour tout $\lambda$, $M\otimes_A B_\lambda$ est un
$B_\lambda$-module \emph{libre}, en vertu de (11.4.1). Mais si l'on pose
$\mathfrak J_\lambda=\mathfrak JB_\lambda$,
$(M\otimes_A B_\lambda)\otimes_{B_\lambda}(B_\lambda/\mathfrak J_\lambda)
=(M\otimes_A(A/\mathfrak J))\otimes_{A/\mathfrak J}
 (B_\lambda/\mathfrak J_\lambda)$
est un $(B_\lambda/\mathfrak J_\lambda)$-module libre, et
$\mathfrak J_\lambda$ est nilpotent. La conclusion résulte donc de l'hypothèse
que $M\otimes_A B_\lambda$ est un $B_\lambda$-module plat et de
$(\mathbf 0_{\mathrm{III}},10.1.2)$.
\end{proof}

\begin{corollary}[11.4.4]
\label{IV.11.4.4-fr}
Soient $A$ un anneau, $M$ un $A$-module~; on suppose qu'il existe un idéal
nilpotent $\mathfrak J$ de $A$ tel que $M\otimes_A(A/\mathfrak J)$ soit un
$(A/\mathfrak J)$-module libre. Alors l'ensemble des idéaux $\mathfrak K$ de
$A$ tels que $M\otimes_A(A/\mathfrak K)$ soit un $(A/\mathfrak K)$-module
libre admet un plus petit élément $\mathfrak K_0$ \textup{(qui est aussi le plus petit
des idéaux $\mathfrak K$ tels que $M\otimes_A(A/\mathfrak K)$ soit un
$(A/\mathfrak K)$-module \emph{plat})}.
\oldpage[IV-3]{146}
Pour qu'un homomorphisme $u:A\to A'$ soit tel que $M\otimes_A A'$ soit un
$A'$-module libre \textup{(resp.} un $A'$-module plat), il faut et il suffit que $u$ se
factorise en $A\to A/\mathfrak K_0\to A'$ \textup{(ou encore que
$\mathfrak K_0A'=0$)}.
\end{corollary}

\begin{proof}
Le fait que l'intersection $\mathfrak K_0$ des idéaux $\mathfrak K$ pour
lesquels $M\otimes_A(A/\mathfrak K)$ est un $(A/\mathfrak K)$-module libre
soit le plus petit de ces idéaux résulte de (11.4.1) appliqué à l'anneau
$A/\mathfrak K_0$, à son idéal nilpotent $\mathfrak J/\mathfrak K_0$ et aux
homomorphismes $A/\mathfrak K_0\to A/\mathfrak K$, dont les noyaux ont $0$
pour intersection. Si $A'$ est une $(A/\mathfrak K_0)$-algèbre, on a
$M\otimes_A A'=
(M\otimes_A(A/\mathfrak K_0))\otimes_{A/\mathfrak K_0}A'$,
donc $M\otimes_A A'$ est un $A'$-module libre. Réciproquement, si $A'$ est
une $A$-algèbre telle que $M\otimes_A A'$ soit un $A'$-module libre et si
$\mathfrak K$ est le noyau de l'homomorphisme $A\to A'$, il résulte de
(11.4.1) appliqué à l'anneau $A/\mathfrak K$, au
$(A/\mathfrak K)$-module $M\otimes_A(A/\mathfrak K)$, à l'idéal nilpotent
$\mathfrak J\mathfrak K/\mathfrak K$ et à l'homomorphisme \emph{injectif}
$A/\mathfrak K\to A'$, que $M\otimes_A(A/\mathfrak K)$ est un
$(A/\mathfrak K)$-module libre, donc que $\mathfrak K\supset\mathfrak K_0$.
Le fait que l'on puisse remplacer «~libre~» par «~plat~» dans ce qui précède
(en conservant naturellement l'hypothèse que $M\otimes_A(A/\mathfrak J)$ est
un $(A/\mathfrak J)$-module \emph{libre}) résulte de
$(\mathbf 0_{\mathrm{III}},10.1.2)$, comme on l'a vu au début de la
démonstration de (11.4.1).
\end{proof}

\begin{proposition}[11.4.5]
\label{IV.11.4.5-fr}
Soient $Y$ un préschéma irréductible, $f:X\to Y$ un morphisme de présentation
finie, $\mathcal F$ un $\mathcal O_X$-Module de présentation finie. Il existe
alors un ouvert non vide $U$ dans $Y$ et un sous-schéma fermé $Z$ de $U$, de
présentation finie sur $U$, ayant la propriété suivante~: pour tout changement
de base $U'\to U$, si l'on pose $X'=X\times_YU'$, $f'=f_{(U')}$ et
$\mathcal F'=\mathcal F\otimes_{\mathcal O_Y}\mathcal O_{U'}$, alors, pour que
$\mathcal F'$ soit $f'$-plat, il faut et il suffit que le morphisme $U'\to U$
se factorise en $U'\to Z\to U$. Un tel schéma $Z$ a même espace sous-jacent que
$U$. Supposons de plus $Y$ affine, et soit $(W_i)$ un recouvrement fini de $X$
par des ouverts affines~; alors, on peut supposer $U$ choisi de sorte que, si
$U'=\operatorname{Spec}(A')$ est un schéma affine et $U'\to U$ un morphisme qui
se factorise en $U'\to Z\to U$, chacun des
$\Gamma(W_i\times_YU',\mathcal F')$ est un $A'$-module libre.
\end{proposition}

\begin{proof}
On peut évidemment se borner au cas où $Y=\operatorname{Spec}(A)$ est affine.
Utilisant (8.9.1), il y a un sous-anneau noethérien $A_1$ de $A$, un morphisme
de type fini $f_1:X_1\to Y_1=\operatorname{Spec}(A_1)$ et un
$\mathcal O_{X_1}$-Module cohérent $\mathcal F_1$ tels que $f=f_{(Y)}$ et
$\mathcal F=\mathcal F_1\otimes_{\mathcal O_{Y_1}}\mathcal O_Y$~; on peut en
outre supposer que les $W_i$ sont images réciproques d'ouverts affines de
$X_1$. Notons en outre que $Y_1$ est irréductible, le morphisme $Y\to Y_1$
étant dominant $(\mathbf I,1.2.7)$. Supposons la proposition démontrée pour
$Y_1$, $X_1$ et $\mathcal F_1$, et soient $U_1$ l'ouvert de $Y_1$ et $Z_1$ le
sous-schéma fermé de $U_1$ ayant les propriétés voulues, $U=U_1\times_{Y_1}Y$
et $Z=Z_1\times_{Y_1}Y$ leurs images réciproques. Alors, si $U'\to U$ est un
changement de base tel que
$\mathcal F'=\mathcal F\otimes_{\mathcal O_Y}\mathcal O_{U'}
=\mathcal F_1\otimes_{\mathcal O_{Y_1}}\mathcal O_{U'}$ soit $f'$-plat, le
morphisme $U'\to U_1$ se factorise en $U'\to Z_1\to U_1$~; mais comme
$U'\to U_1$ se factorise aussi en $U'\to U\to U_1$, la définition du produit
de préschémas montre que $U'\to U$ se factorise en $U'\to Z\to U$.

On peut donc se borner au cas où $A$ est \emph{noethérien}~; soit
$\mathfrak N$ son nilradical, qui est nilpotent, et posons
$A_0=A_{\mathrm{red}}=A/\mathfrak N$,
$Y_0=\operatorname{Spec}(A_0)=Y_{\mathrm{red}}$,
$X_0=X\times_YY_0$, $\mathcal F_0=\mathcal F\otimes_{\mathcal O_Y}
\mathcal O_{Y_0}$, $W_{i0}=W_i\times_YY_0$~; si
$M_i=\Gamma(W_i,\mathcal F)$, $M_{i0}=\Gamma(W_{i0},\mathcal F_0)$ est donc
égal à $M_i\otimes_A A_0$. Comme $A_0$ est \emph{intègre}, on peut, en vertu du
théorème de platitude générique (6.9.2), en remplaçant au besoin $Y$ par un
ouvert non vide de $Y$, supposer que les $M_{i0}$ soient des $A_0$-modules
\emph{libres}. En vertu de (11.4.4), il y a donc pour chaque $i$ un idéal
$\mathfrak J_i\subset\mathfrak N$ tel
\oldpage[IV-3]{147}
que les $A$-algèbres $A'$ pour lesquelles les $M_i\otimes_A A'$ sont des
$A'$-modules libres (ou plats) sont exactement celles pour lesquelles
$\mathfrak J_iA'=0$. Soit $\mathfrak J=\sum_i\mathfrak J_i$, qui est un idéal
contenu dans $\mathfrak N$~; dire que $\mathcal F\otimes_A A'$ est $A'$-plat
équivaut à dire que les $M_i\otimes_A A'$ sont tous des $A'$-modules plats,
donc que $\mathfrak JA'=0$. Il en résulte que si l'on prend
$Z=V(\mathfrak J)$, on répond à la question, car pour qu'un morphisme
$U'\to U$ ait la propriété de l'énoncé, il faut et il suffit évidemment que
pour tout ouvert affine $V'\subset U'$, le morphisme $V'\to U$ ait la même
propriété.
\end{proof}

\begin{corollary}[11.4.6]
\label{IV.11.4.6-fr}
Soient $A$ un anneau tel que $\operatorname{Spec}(A)$ soit irréductible,
$\mathfrak p$ son unique idéal premier minimal, $B$ une $A$-algèbre de
présentation finie, $M$ un $B$-module de présentation finie. Supposons que
$M_\mathfrak p$ soit un $A_\mathfrak p$-module plat~; alors il existe
$t\in A\setmin\mathfrak p$ tel que $M_t$ soit un $A_t$-module libre.
\end{corollary}

\begin{proof}
Appliquant (11.4.5) à $Y=\operatorname{Spec}(A)$, $X=\operatorname{Spec}(B)$,
$\mathcal F=\widetilde M$, on peut (en remplaçant au besoin $A$ par $A_h$, où
$h\in A\setmin\mathfrak p$) supposer qu'il existe un idéal de type fini
$\mathfrak J$ dans $A$ tel que les $A$-algèbres $A'$ pour lesquelles
$M\otimes_A A'$ soit un $A'$-module libre (ou plat) sont exactement celles
telles que $\mathfrak JA'=0$. En particulier, l'hypothèse que
$M_\mathfrak p$ soit un $A_\mathfrak p$-module plat entraîne
$\mathfrak JA_\mathfrak p=0$, ou encore $\mathfrak J_\mathfrak p=0$. Mais
comme $\mathfrak J$ est de type fini, il existe $t\in A\setmin\mathfrak p$ tel que
$\mathfrak J_t=0$, ou encore $\mathfrak JA_t=0$, et par suite $M_t$ est un
$A_t$-module libre.
\end{proof}

\begin{proposition}[11.4.7]
\label{IV.11.4.7-fr}
Soient $A$ un anneau noethérien, $\mathfrak J$ un idéal de $A$, $M$ un
$A$-module idéalement séparé $(\mathbf 0_{\mathrm{III}},10.2.1)$ pour la
topologie $\mathfrak J$-préadique. Soit $(\mathfrak p_i)_{1\leq i\leq r}$ une
famille finie d'idéaux premiers de $A$ contenant $\mathfrak J$~; pour tout
entier $n\geq0$, soit $\mathfrak p_i^{(n)}$ la puissance symbolique $n$-ième de
$\mathfrak p_i$ (noyau de l'homomorphisme
$A\to A_{\mathfrak p_i}/\mathfrak p_i^nA_{\mathfrak p_i}$)~; posons
$\mathfrak J_n=\bigcap_{i=1}^{r}\mathfrak p_i^{(n)}$, de sorte que
$\mathfrak J_n\supset\mathfrak J^n$, et supposons que la topologie définie par
la filtration $(\mathfrak J_n)$ soit identique à la topologie
$\mathfrak J$-préadique (autrement dit que, pour tout $n$, il existe $m$ tel
que $\mathfrak J^n\supset\mathfrak J_m$). Pour que $M$ soit un $A$-module
plat, il faut et il suffit que $M\otimes_A(A/\mathfrak J)$ soit un
$(A/\mathfrak J)$-module plat et que, pour tout $i$,
$M\otimes_AA_{\mathfrak p_i}=M_{\mathfrak p_i}$ soit un
$A_{\mathfrak p_i}$-module plat.
\end{proposition}

\begin{proof}
Comme $M$ est idéalement séparé, il suffit, en vertu de
$(\mathbf 0_{\mathrm{III}},10.2.1)$, de montrer que, pour tout $n\geq0$,
$M\otimes_A(A/\mathfrak J^n)$ est un $(A/\mathfrak J^n)$-module plat~;
puisque tout $\mathfrak J^n$ contient un $\mathfrak J_m$, il revient au même
de prouver que pour tout $n$, $M\otimes_A(A/\mathfrak J_n)$ est un
$(A/\mathfrak J_n)$-module plat. Or, comme
$\mathfrak J_1\supset\mathfrak J$, $M\otimes_A(A/\mathfrak J_1)$ est un
$(A/\mathfrak J_1)$-module plat~; dans l'anneau $A/\mathfrak J_n$, l'idéal
$\mathfrak J_1/\mathfrak J_n$ est nilpotent et enfin l'intersection des noyaux
des homomorphismes
$A/\mathfrak J_n\longrightarrow
A_{\mathfrak p_i}/\mathfrak p_i^nA_{\mathfrak p_i}$
est nulle dans $A/\mathfrak J_n$, par définition de $\mathfrak J_n$. Il suffit
donc, par (11.4.1), de vérifier que
$M\otimes_A(A_{\mathfrak p_i}/\mathfrak p_i^nA_{\mathfrak p_i})$
est un $(A_{\mathfrak p_i}/\mathfrak p_i^nA_{\mathfrak p_i})$-module plat, ce
qui résulte de l'hypothèse que $M\otimes_AA_{\mathfrak p_i}$ est un
$A_{\mathfrak p_i}$-module plat.
\end{proof}

\begin{remark}[11.4.8]
\label{IV.11.4.8-fr}
L'hypothèse faite dans (11.4.7) sur la topologie définie par les
$\mathfrak J_n$ est vérifiée si, pour tout $n$ assez grand,
$\operatorname{Ass}(A/\mathfrak J^n)$ \emph{est contenu dans l'ensemble des}
$\mathfrak p_i$. En effet, $\mathfrak J^n$ est alors intersection d'idéaux
primaires pour les $\mathfrak p_i$, dont chacun contient une puissance
symbolique de $\mathfrak p_i$, d'où la conclusion. En particulier~:
\end{remark}

\begin{corollary}[11.4.9]
\label{IV.11.4.9-fr}
Soient $A$ un anneau noethérien, $\mathfrak J$ un idéal nilpotent de $A$, $M$
un $A$-module. Pour que $M$ soit un $A$-module plat \textup{(resp.} libre), il faut et
il suffit que $M\otimes_A(A/\mathfrak J)$ soit un $(A/\mathfrak J)$-module
plat \textup{(resp.} libre) et que pour tout idéal premier
$\mathfrak p\in\operatorname{Ass}(A)$, $M_\mathfrak p$ soit un
$A_\mathfrak p$-module plat.
\end{corollary}
% IV3_B22_APPEND_ANCHOR

\oldpage[IV-3]{148}

L'assertion relative au cas où $M\otimes_A(A/\mathfrak J)$ est libre se déduit
encore de l'assertion relative au cas où $M\otimes_A(A/\mathfrak J)$ est plat
par $(\mathbf 0_{\mathrm{III}},10.1.2)$.

\begin{corollary}[11.4.10]
\label{IV.11.4.10-fr}
Soient $A$ un anneau noethérien, $\mathfrak J$ un idéal de $A$, $M$ un
$A$-module. On suppose que $M$ est idéalement séparé pour la topologie
$\mathfrak J$-préadique et que $\operatorname{gr}_{\mathfrak J}(A)$ est un
$(A/\mathfrak J)$-module plat. Pour que $M$ soit un $A$-module plat, il faut et
il suffit que $M\otimes_A(A/\mathfrak J)$ soit un $(A/\mathfrak J)$-module plat
et que pour tout $\mathfrak p\in\operatorname{Ass}(A/\mathfrak J)$,
$M_{\mathfrak p}$ soit un $A_{\mathfrak p}$-module plat.
\end{corollary}

\begin{proof}
Compte tenu de (11.4.8), tout revient à montrer que
$\operatorname{Ass}(A/\mathfrak J^n)$ est contenu dans
$\operatorname{Ass}(A/\mathfrak J)$ pour tout $n$. Or, si $a\in A$ n'appartient
à aucun des $\mathfrak p\in\operatorname{Ass}(A/\mathfrak J)$, l'homothétie de
rapport $a$ est injective dans $A/\mathfrak J$~; comme chacun des
$\mathfrak J^k/\mathfrak J^{k+1}$ est un $(A/\mathfrak J)$-module plat, $a$ est
aussi un élément $(\mathfrak J^k/\mathfrak J^{k+1})$-régulier, donc $a$ est
$(A/\mathfrak J^n)$-régulier pour tout $n$ (Bourbaki, \emph{Alg. comm.},
chap.~III, \S~2, n\textsuperscript{o}~8, cor.~1 du th.~1), et par suite n'appartient à
aucun idéal premier associé à $A/\mathfrak J^n$, d'où le corollaire (Bourbaki,
\emph{Alg. comm.}, chap.~II, \S~1, n\textsuperscript{o}~1, prop.~2).
\end{proof}

\begin{proposition}[11.4.11]
\label{IV.11.4.11-fr}
Soient $A$ un anneau local artinien, d'idéal maximal $\mathfrak m$,
$Y=\operatorname{Spec}(A)$, $y$ l'unique point de $Y$, $f:X\to Y$ un morphisme
localement de type fini, $\mathcal F$ un $\mathcal O_X$-Module cohérent. Soient
$(A'_\alpha)$ une famille d'anneaux locaux, et pour chaque $\alpha$,
$u_\alpha:A\to A'_\alpha$ un homomorphisme d'anneaux (nécessairement local).
Posons $Y'_\alpha=\operatorname{Spec}(A'_\alpha)$,
$X'_\alpha=X\times_Y Y'_\alpha$,
$\mathcal F'_\alpha=\mathcal F\otimes_A A'_\alpha$. Soit $x$ un point de $X$,
et supposcns vérifiées les conditions suivantes~:
\begin{enumerate}[label=\textup{(\roman*)}]
  \item L'intersection des noyaux des $u_\alpha$ est réduite à $0$.
  \item L'extension $\boldsymbol{k}(x)$ du corps résiduel $k$ de $A$ est primaire \textup{(4.3.1)}
  (condition automatiquement vérifiée si $k$ est séparablement clos).
  \item Pour tout $\alpha$, il existe un point $x'_\alpha\in X'_\alpha$ dont
  les projections respectives dans $X$ et $Y'_\alpha$ sont $x$ et le point
  fermé $y'_\alpha$ de $Y'_\alpha$, et tel que $\mathcal F'_\alpha$ soit
  $Y'_\alpha$-plat au point $x'_\alpha$.
\end{enumerate}

Dans ces conditions, $\mathcal F$ est $f$-plat au point $x$.
\end{proposition}

\begin{proof}
La question étant locale sur $X$, on peut évidemment se borner au cas où $f$
est un morphisme de type fini et supposer $\mathcal F_x\neq0$. Nous procéderons
en plusieurs étapes.

\medskip
\noindent I) \emph{Réduction au cas où $A'_\alpha$ est un anneau local
d'un $Y$-préschéma de type fini et où le corps résiduel $k'_\alpha$ de
$A'_\alpha$ est une extension finie de $k$.}

\medskip
Comme la réduction se fait séparément sur chaque $A'_\alpha$, on peut supprimer
dans cette partie l'indice $\alpha$. Soit $\mathfrak m'$ l'idéal maximal de
$A'$. Considérons $A'$ comme limite inductive de ses sous-$A$-algèbres de type
fini $B_\lambda$, et posons $\mathfrak n_\lambda=\mathfrak m'\cap B_\lambda$~;
$A'$ est aussi limite inductive de ses sous-anneaux \emph{locaux}
$B'_\lambda=(B_\lambda)_{\mathfrak n_\lambda}$ (5.13.3), et on a évidemment
$\mathfrak m'\cap B'_\lambda=\mathfrak n'_\lambda$, idéal maximal de
$B'_\lambda$. Il existe donc (11.2.6) un indice $\lambda$ tel que, en posant
$Z'_\lambda=\operatorname{Spec}(B'_\lambda)$,
$\mathcal F\otimes_A B'_\lambda$ soit $Z'_\lambda$-plat au point $x'_\lambda$,
projection de $x'$, la projection de $x'_\lambda$ sur $Z'_\lambda$ étant le
\emph{point fermé} $z'_\lambda$ de $Z'_\lambda$.

On peut donc supposer qu'il existe un schéma affine $Z'$ de type fini sur $Y$
et un point $z'$ de $Z'$ tel que $A'=\mathcal O_{Z',z'}$, et que, si l'on pose
$T'=X\times_Y Z'$, il existe un point $\boldsymbol{t}'\in T'$ dont les projections dans
$Z'$ et $X$ sont $z'$ et $x$, et tel que $\mathcal F\otimes_Y Z'$ soit
$Z'$-plat au point $t'$. Soient $Z'_1$ (resp. $X_1$) un sous-préschéma fermé de
$Z'$ (resp. $X$) ayant pour espace sous-jacent l'adhérence de $\{z'\}$ (resp.
$\{x\}$), et posons $T'_1=X_1\times_Y Z'_1$. L'ensemble $U$ des points de
$T'$ où $\mathcal F\otimes_Y Z'$ est $Z'$-plat est ouvert dans $T'$ (11.1.1) et
contient $t'$, donc
\oldpage[IV-3]{149}
$V=U\cap T'_1$ est ouvert non vide dans $T'_1$. L'anneau $A$, étant artinien,
est un anneau de Jacobson, donc (10.4.6) $Z'_1$ et $T'_1$ sont des préschémas de
Jacobson~; il existe par suite dans $V$ un point $t'_0$ fermé dans $V$, et son
image $z'_0$ dans $Z'_1$ est un point \emph{fermé} de $Z'_1$ (10.4.7). Soient $f_1$ la
restriction de $f$ à $X_1$, $f'_1:T'_1\to Z'_1$ la projection canonique~;
$V\cap f_1'^{-1}(z'_0)$ est un ouvert non vide dans
$f_1^{-1}(y)\otimes_{\mathbf{k}(y)}\boldsymbol{k}(z'_0)$, et comme ce dernier préschéma est plat sur
$f_1^{-1}(y)$, un point maximal $t'_1$ de $V\cap f_1'^{-1}(z'_0)$ est
nécessairement au-dessus de l'unique point maximal $x$ de $X_1$ (2.3.4).
Enfin, $\boldsymbol{k}(z'_0)$ est une extension \emph{finie} de $k$ (\textbf I, 6.4.2) et
l'homomorphisme $A=\mathcal O_{Y,y}\to A'=\mathcal O_{Z',z'}$ se factorise en
$\mathcal O_{Y,y}\to\mathcal O_{Z',z'_0}\to\mathcal O_{Z',z'}$, donc son noyau
est contenu dans celui de $\mathcal O_{Y,y}\to\mathcal O_{Z',z'_0}$. Ceci
achève la réduction annoncée.

\medskip
\noindent II) \emph{Réduction au cas où les $A'_\alpha$ sont en nombre
fini et sont des $A$-algèbres finies.} --- Soit $\mathfrak m'_\alpha$ l'idéal
maximal de $A'_\alpha$~; comme $A'_\alpha$ est un anneau local noethérien,
l'intersection des $\mathfrak m'^n_\alpha$ est $0$ $(\mathbf 0_{\mathrm I},7.3.5)$~;
l'intersection des $u_\alpha^{-1}(\mathfrak m'^n_\alpha)$, pour tous les indices
$n$ et $\alpha$, est donc égale à l'intersection des noyaux des $u_\alpha$, donc
est réduite à $0$ par l'hypothèse (i). Puisque $A$ est artinien, il y a un nombre
\emph{fini} de ces idéaux dont l'intersection est déjà $0$~; notons-les
$u_i^{-1}(\mathfrak m_i'^{\,n_i})$ $(1\leq i\leq r)$. Comme les $A'_i$ sont
noethériens, les $\mathfrak m_i'^{\,h}/\mathfrak m_i'^{\,h+1}$ sont des
$(A'_i/\mathfrak m'_i)$-modules de longueur finie, et comme
$A'_i/\mathfrak m'_i$ est un $k$-espace vectoriel de rang fini, on voit que
$A''_i=A'_i/\mathfrak m_i'^{\,n_i}$ est une $A$-algèbre \emph{finie} et un anneau local
artinien. La réduction annoncée résulte donc de (2.1.4).

\medskip
\noindent III) \emph{Fin de la démonstration.} --- Supposons désormais
que les $A'_i$ $(1\leq i\leq r)$ soient en nombre fini et soient des
$A$-algèbres finies. Pour tout $i$, le corps résiduel $k_i$ de $A'_i$ est une
extension finie de $k$~; utilisant l'hypothèse (ii), on en conclut que l'image
réciproque de $x$ dans $X'_i$ est réduite \emph{au seul point} $x'_i$ (4.3.2). Soient
$Y'$ le préschéma \emph{somme} des $Y'_i$, $X'=X\times_Y Y'$, somme des $X'_i$,
$\mathcal F'=\mathcal F\otimes_Y Y'$. L'hypothèse entraîne que
$\mathcal F'$ est $Y'$-plat aux points de l'image réciproque de $x$ par la
projection $p:X'\to X$. Comme $p$ est de type fini, il existe par suite un
\emph{ouvert} $U'\supset p^{-1}(x)$ tel que $\mathcal F'$ soit $Y'$-plat aux points de
$U'$ (11.1.1). En outre, le morphisme $Y'\to Y$ est \emph{fini} puisque les $A'_i$
sont des $A$-algèbres finies~; donc $p$ est un morphisme fini (\textbf{II},
6.1.5), par suite \emph{fermé} (\textbf{II}, 6.1.10), et il existe donc un voisinage
affine $U$ de $x$ dans $X$ tel que $p^{-1}(U)\subset U'$. Soient $B$ et $A'$
les anneaux des schémas $U$ et $Y'$ ($A'$ étant le composé direct des $A'_i$)~;
remplaçant $X$ par $U$, on a donc $\mathcal F=\widetilde M$, où $M$ est un
$B$-module, et par hypothèse $M'=M\otimes_A A'$ est un $A'$-module \emph{plat}
(2.1.2)~; comme l'homomorphisme $A\to A'$ est injectif par construction, on
peut appliquer (11.4.3), qui prouve que $M$ est un $A$-module plat, C.Q.F.D.
\end{proof}

\begin{proposition}[11.4.12]
\label{IV.11.4.12-fr}
Soient $A$ un anneau local artinien de corps résiduel $k$,
$Y=\operatorname{Spec}(A)$, $f:X\to Y$ un morphisme localement de type fini,
$\mathcal F$ un $\mathcal O_X$-Module cohérent. Soient $(A'_\alpha)$ une famille
d'anneaux locaux, et pour chaque $\alpha$, $u_\alpha:A\to A'_\alpha$ un
homomorphisme d'anneaux. Posons $Y'_\alpha=\operatorname{Spec}(A'_\alpha)$,
$X'_\alpha=X\times_Y Y'_\alpha$, $f'_\alpha=f_{(Y'_\alpha)}$,
$\mathcal F'_\alpha=\mathcal F\otimes_A A'_\alpha$. Soit $x$ un point de $X$,
et supposons vérifiées les conditions suivantes~:
\begin{enumerate}[label=\textup{(\roman*)}]
  \item L'intersection des noyaux des $u_\alpha$ est réduite à $0$.
  \item Pour tout $\alpha$, $\mathcal F'_\alpha$ est $f'_\alpha$-plat en \textup{tous}
  les points $x'_\alpha\in X'_\alpha$ dont les projections respectives dans
  $X$ et $Y'$ sont $x$ et le point fermé $y'_\alpha$ de $Y'_\alpha$.
\end{enumerate}

Alors $\mathcal F$ est $f$-plat au point $x$.
\end{proposition}

% IV3_B22_P149_ANCHOR

\oldpage[IV-3]{150}
\begin{proof}
Par hypothèse, l'intersection des noyaux des $u_\alpha$ est réduite à $0$~;
comme $A$ est artinien, il y a déjà un nombre \emph{fini} de ces noyaux dont
l'intersection est $0$, donc on peut se borner au cas où la famille
$(A'_\alpha)$ est \emph{finie}. Soit $k'$ une clôture algébrique de $k$~; on sait
$(\mathbf 0_{\mathrm{III}},10.3.1)$ qu'il existe un homomorphisme local
$A\to B$ faisant de $B$ un $A$-module plat, tel que $B$ soit un anneau local
\emph{artinien}, \emph{entier} sur $A$ et que $B\otimes_A k$ soit isomorphe à $k'$.
Par platitude, les noyaux des homomorphismes
$B\to B'_\alpha=B\otimes_A A'_\alpha$ se déduisent de ceux des $u_\alpha$ par
tensorisation avec $B$, et comme ils sont en nombre fini, leur intersection est
réduite à $0$ $(\mathbf 0_{\mathrm I},6.1.3)$. Considérons les anneaux
$B'_{\alpha\beta}$, localisés de $B'_\alpha$ en ses idéaux \emph{maximaux}~; on sait
(Bourbaki, \emph{Alg. comm.}, chap.~II, \S~3, n\textsuperscript{o}~3, cor.~2 du th.~1)
que l'intersection des noyaux des homomorphismes
$B'_\alpha\to B'_{\alpha\beta}$ (pour un $\alpha$ donné) est réduite à $0$~;
on en conclut que l'intersection des noyaux des homomorphismes composés
$v_{\alpha\beta}:B\to B'_\alpha\to B'_{\alpha\beta}$ ($\alpha$ et $\beta$
variables) est réduite à $0$. D'autre part, comme $B$ est entier sur $A$,
$B'_\alpha$ est entier sur $A'_\alpha$, donc ses idéaux maximaux sont au-dessus
de l'idéal maximal de $A'_\alpha$. Si l'on pose
$Z'_{\alpha\beta}=\operatorname{Spec}(B'_{\alpha\beta})$,
$T'_{\alpha\beta}=X'_\alpha\times_{Y'_\alpha}Z'_{\alpha\beta}$,
$f'_{\alpha\beta}=f_{(Z'_{\alpha\beta})}$,
$\mathcal G=\mathcal F\otimes_A B$ et
$\mathcal G'_{\alpha\beta}=\mathcal G\otimes_B B'_{\alpha\beta}$, on voit donc
que les hypothèses (i) et (ii) sont encore vérifiées lorsqu'on remplace
respectivement $A$, $A'_\alpha$, $u_\alpha$, $\mathcal F$ et $x$ par $B$,
$B'_{\alpha\beta}$, $v_{\alpha\beta}$, $\mathcal G$ et un point $t$ de
$T=X\otimes_A B$ au-dessus de $x$. Comme le corps résiduel de $B$ est
séparablement clos, on déduit de (11.4.11) que $\mathcal G$ est plat sur $B$ au
point $t$. Mais puisque $B$ est un $A$-module fidèlement plat, on conclut de
(2.1.4) que $\mathcal F$ est plat sur $A$ au point $x$, ce qui prouve (11.4.12).
\end{proof}

\begin{corollary}[11.4.13]
\label{IV.11.4.13-fr}
Soient $A$ un anneau artinien local, $Y=\operatorname{Spec}(A)$,
$f:X\to Y$ un morphisme localement de type fini, $\mathcal F$ un
$\mathcal O_X$-Module cohérent. Soit $(Y'_\alpha)$ une famille de
$A$-préschémas, et pour tout $\alpha$, posons
$X'_\alpha=X\times_Y Y'_\alpha$, $f'_\alpha=f_{(Y'_\alpha)}$,
$\mathcal F'_\alpha=\mathcal F\otimes_Y Y'_\alpha$. Soit $x$ un point de $X$ et
supposons vérifiées les hypothèses suivantes~:
\begin{enumerate}[label=\textup{(\roman*)}]
  \item L'intersection des noyaux des homomorphismes
  $A\to\Gamma(Y'_\alpha,\mathcal O_{Y'_\alpha})$ correspondant aux morphismes
  structuraux est réduite à $0$.
  \item Pour tout $\alpha$, $\mathcal F'_\alpha$ est $f'_\alpha$-plat en \textup{tous}
  les points $x'_\alpha\in X'_\alpha$ dont la projection dans $X$ est $x$.
\end{enumerate}

Alors $\mathcal F$ est $f$-plat au point $x$.
\end{corollary}

\begin{proof}
En effet, pour tout $y'_\alpha\in Y'_\alpha$, considérons le schéma local
$Y''_{\alpha,y'_\alpha}=\operatorname{Spec}(\mathcal O_{y'_\alpha})$~;
en vertu de (2.1.4),
$\mathcal F'_\alpha\otimes_{Y'_\alpha}Y''_{\alpha,y'_\alpha}$ est
$Y''_{\alpha,y'_\alpha}$-\emph{plat} aux points dont les projections sur $X$ et
$Y''_{\alpha,y'_\alpha}$ sont $x$ et le point fermé de
$Y''_{\alpha,y'_\alpha}$. En outre le noyau de l'homomorphisme
$A\to\Gamma(Y'_\alpha,\mathcal O_{Y'_\alpha})$ est l'intersection des noyaux
des homomorphismes
$A\to\Gamma(Y''_{\alpha,y'_\alpha},\mathcal O_{Y''_{\alpha,y'_\alpha}})$, car
on se ramène aussitôt au cas où $Y'_\alpha$ est affine, et il suffit alors
d'appliquer Bourbaki, \emph{Alg. comm.}, chap.~II, \S~3, n\textsuperscript{o}~3, cor.~2 du
th.~1. En remplaçant la famille $(Y'_\alpha)$ par la famille des
$Y''_{\alpha,y'_\alpha}$, on est donc ramené à (11.4.12).
\end{proof}

\subsection{Descente de la platitude par des morphismes quelconques : cas général.}
\label{subsection:IV.11.5-fr}

\begin{theorem}[11.5.1]
\label{IV.11.5.1-fr}
Soient $Y$ un préschéma localement noethérien, $f:X\to Y$ un morphisme
localement de type fini, $\mathcal F$ un $\mathcal O_X$-Module cohérent, $x$ un
point de $X$, $y=f(x)$. Soit $(Y'_\alpha)$ une famille de $Y$-préschémas locaux
$Y'_\alpha=\operatorname{Spec}(A'_\alpha)$ telle que les images des points
fermés $y'_\alpha$ des $Y'_\alpha$ soient toutes égales à $y$. Pour tout
$\alpha$, soient $\mathfrak m'_\alpha$ l'idéal maximal de $A'_\alpha$, et
$u_\alpha:\mathcal O_y\to A'_\alpha$
\oldpage[IV-3]{151}
% IV3_B22_P151_ANCHOR
l'homomorphisme canonique
\textup{(\textbf I, 2.4.4)}~; on suppose que les intersections finies des idéaux
$u_\alpha^{-1}(\mathfrak m_\alpha'^{\,n_\alpha})$ forment un système fondamental
de voisinages de $0$ dans $\mathcal O_y$. Posons
$X'_\alpha=X\times_Y Y'_\alpha$, $f'_\alpha=f_{(Y'_\alpha)}$,
$\mathcal F'_\alpha=\mathcal F\otimes_Y Y'_\alpha$, et supp